\def\stat{kovalev-tarasov}





\def\tit{АРХИТЕКТУРА ПЛАТФОРМЫ ДЛЯ ПРОВЕДЕНИЯ ДВИЖИМЫХ ГИПОТЕЗАМИ ВИРТУАЛЬНЫХ 


ЭКСПЕРИМЕНТОВ$^*$}





\def\titkol{Архитектура платформы для проведения %движимых гипотезами 


виртуальных 


экспериментов}





\def\aut{Д.\,Ю.~Ковалев$^1$, Е.\,А.~Тарасов$^2$, В.\,Н.~Захаров$^3$, 


Н.\,М.~Филимонов$^4$}





\def\autkol{Д.\,Ю.~Ковалев, Е.\,А.~Тарасов, В.\,Н.~Захаров, Н.\,М.~Филимонов}





\titel{\tit}{\aut}{\autkol}{\titkol}





\index{Ковалев Д.\,Ю.}


\index{Тарасов Е.\,А.}


\index{Захаров В.\,Н.}


\index{Филимонов Н.\,М.}


\index{Kovalev D.\,Y.}


\index{Tarasov E.\,A.}


\index{Zakharov V.\,N.}


\index{Filimonov N.\,M.}





{\renewcommand{\thefootnote}{\fnsymbol{footnote}}


\footnotetext[1]


{Работа выполнена при частичной финансовой поддержке РФФИ (проект 18-07-01434~А).}} 





\renewcommand{\thefootnote}{\arabic{footnote}}


\footnotetext[1]{Институт проблем информатики Федерального исследовательского 


центра <<Информатика и~управление>> Российской академии наук, 


\mbox{dkovalev@ipiran.ru}}


\footnotetext[2]{Институт проблем информатики Федерального исследовательского 


центра <<Информатика и~управление>> Российской академии наук, 


\mbox{etarasov@outlook.com}}


\footnotetext[3]{Институт проблем информатики Федерального исследовательского 


центра <<Информатика и~управление>> Российской академии наук, 


\mbox{vzakharov@ipiran.ru}}


\footnotetext[4]{Институт информационных и~вычислительных технологий Национального 


исследовательского института <<Московский энергетический институт>>, 


\mbox{filimonovn160@gmail.com}}





 \label{st\stat}





 \thispagestyle{headings}


 


 \vspace*{-12pt}





\Abst{Рассматривается проблема проведения движимых гипотезами виртуальных 


экспериментов. Проведен анализ методов управления виртуальными экспериментами 


в~существующих системах работы с~экспериментами, по результатам которого 


сформирован жизненный цикл виртуального эксперимента. Для этапов жизненного цикла 


представлены основные операции управления виртуальными экспериментами, а также 


основные этапы по работе эксперта с~платформой исполнения виртуального 


эксперимента. Предложена программная архитектура платформы управления 


виртуальными экспериментами и~гипотезами с~описанием основных компонентов 


платформы и~их функций, которые реализуют основные операции жизненного цикла.}





\KW{жизненный цикл виртуального эксперимента; гипотеза; интенсивное использование 


данных}





\DOI{10.14357\!/\!08696527200206} 





\vskip 10pt plus 9pt minus 6pt





 \thispagestyle{headings}


 


 \vspace*{-12pt} 





\section{Введение}





В настоящее время в~научных исследованиях необходимо решать задачи по 


получению новых знаний на основе разноструктурированных распределенных 


данных. Наборы данных получают в~результате наблюдений, экспериментов или 


компьютерного моделирования, а также построения математических моделей для 


интерпретации наблюдаемых явлений. Такие исследования опираются на явное 


использование гипотез и~реализующих их вычислительных моделей, 


используемых для преодоления разрыва между концептуальным описанием 


изучаемого явления и~его симуляцией. 





Проект 18-07-01434, выполняемый Федеральным научным центром 


<<Информатика и~управление>> Российской академии наук при финансовой 


поддержке Российского фонда фундаментальных исследований, направлен на 


разработку методов и~средств проведения виртуальных экспериментов, 


основанных на проверке гипотез, в~исследованиях с~интенсивным 


использованием данных для их ускорения. 





Виртуальный эксперимент в~рамках исследований определяется как кортеж 


$\langle O, H, M, R, W, C\rangle$~\cite{1-kt}, где~$O$~--- это онтология 


предметной области, которая представляет собой набор понятий и~отношений, 


определенный на некотором формальном языке (например, OWL~2~\cite{2-kt}); 


$H$~--- набор спецификаций гипотез и~взаимосвязей между ними (описание 


гипотез в~работе ограничивается системами уравнений), являющийся частью 


онтологии~$O$ и~использующий ее понятия; $M$~--- набор моделей, в~котором 


каждая модель представляет собой набор программных функций, реализующих 


спецификацию некоторой гипотезы; $R: M \to H$~--- сюръективное отображение 


из множества моделей в~множество гипотез, соответствующее реализации 


гипотез моделями; $W$~--- поток работ, т.\,е.\ набор задач, организованный 


с~помощью определенных конструкций (образцов потоков работ, таких как 


разбиение, объединение и~т.\,д.\cite{3-kt}) и~реализующий эксперимент (каж\-дая 


задача представляет собой вызов функций некоторой модели из~$M$); $C$~--- 


конфигурация эксперимента, представляющая собой отображение задач потока 


работ в~набор значений параметров функций. При проведении нового 


виртуального эксперимента требуется определить полностью кортеж $\langle O, 


H, M, R, W, C\rangle$. Пример определения элементов кортежа виртуального 


эксперимента для задачи моделирования масс звезд Безансонской модели 


Галактики приведен в~работе~\cite{1-kt}.





Анализ возможностей и~ограничений существующих систем управления 


виртуальными экспериментами показывает необходимость разработки 


и~реализации методов и~инструментов, имеющих дело с~наборами зависимых 


гипотез, называемых далее \textit{решетками гипотез}. В~работе~\cite{4-kt} 


приведен алгоритм построения решеток гипотез для виртуального 


эксперимента, а также дана оценка сложности алгоритма. Построение решеток 


гипотез вносит ряд улучшений в~проведение виртуального эксперимента, 


например: использование промежуточных результатов вычисления независимых 


гипотез приводит к~ускорению вычислений. Порождение гипотез представляет 


дополнительную трудность, а потому при не\-об\-хо\-ди\-мости гипотезы могут 


автоматически генерироваться из данных. Так, в~работе~\cite{5-kt} рассмотрены 


подходы к~извлечению гипотез в~виде нелинейных функциональных 


зависимостей между несколькими переменными.





Данная статья посвящена разработке архитектуры платформы управления 


виртуальными экспериментами, включающей средства концептуального 


пред\-став\-ле\-ния гипотез в~процессе проведения виртуальных экспериментов, 


сохранения эволюции гипотез и~моделей, построения связей между несколькими 


гипотезами в~одном эксперименте, воспроизводимости экспериментов. 


Платформа предполагает исследование взаимосвязи между гипотезами в~форме 


математических уравнений, предоставляя информацию не только о зависимостях 


переменных в~одной гипотезе\!/\!модели, но также и~о нескольких зависимых 


гипотезах в~одном виртуальном эксперименте. Одна из важнейших целей 


платформы~--- повышение скорости проведения исследований за счет повторного 


использования уже вычисленных фрагментов виртуальных экспериментов.





Работа структурирована следующим образом. В~разд.~2 приводится анализ 


существующих методов управления экспериментами, по результатам которого 


выделяются жизненный цикл и~основные операции формального 


манипулирования виртуальными экспериментами и~гипотезами. В~разд.~3 


представленa архитектура платформы для проведения движимых гипотезами 


виртуальных экспериментов, реализующих операции всех этапов жизненного 


цикла эксперимента. 





\section{Методы и~операции управления виртуальными 


экспериментами и~гипотезами}





В данном разделе рассматриваются методы манипулирования гипотезами 


и~экспериментами в~ряде существующих систем работы с~виртуальными 


экспериментами. По результатам анализа методов манипулирования 


виртуальными экспериментами выделены операции для последующей реализации 


в~рамках платформы управления виртуальными экспериментами.





В системе Hephaestus\cite{6-kt} основными методами для манипулирования 


гипотезами и~виртуальными экспериментами являются: 


\begin{enumerate}[(1)]


\item метод 


автоматического по\-рож\-де\-ния гипотез из данных;


\item метод интервенций при 


исполнении виртуального эксперимента;


\item метод корректного разбиения данных 


наблюдений при проверке статистических гипотез;


\item метод ранжирования 


нескольких конкурирующих гипотез;


\item метод построения графа 


при\-чин\-но-след\-ст\-вен\-ных и~вероятностных зависимостей.


\end{enumerate}





 В~качестве \textit{метода 


автоматического порождения гипотез} используется символьная 


регрессия~\cite{7-kt} при подборе аппроксимирующей формулы для выбранного 


подмножества переменных на наборе данных. Сущность \textit{метода 


интервенций} при исполнении виртуального эксперимента заключается 


в~уста\-нов\-ле\-нии существующих и~возможной генерации новых гипотез, которые 


значимо влияют на выбранную экспертом гипотезу~\cite{8-kt}. \textit{Метод 


корректного разбиения данных наблюдений} при проверке статистических 


гипотез требуется для обеспечения несмещенных оценок при дальнейшем 


тестировании гипотез~\cite{9-kt}. \textit{Метод ранжирования} нескольких 


конкурирующих гипотез позволяет по вы\-бран\-ной мет\-ри\-ке (например, 


достигаемому уровню значимости) или нескольким мет\-ри\-кам выбрать одну или 


несколько гипотез, наиболее соответствующих наблюдениям. \textit{Постро\-ение 


графа причинно-следственных и~вероятностных зависимостей} применяется для 


оценки соответствия полученных зависимостей общепринятым в~данной области, 


а также для потенциального исследования значимых корреляций и~их 


последующего преобразования в~при\-чин\-но-след\-ст\-вен\-ные зависимости. 


Кроме того, построенные графы зависимостей могут в~явном виде передаваться 


другим экспертам в~данной об\-ласти.





В системе Upsilon-DB~\cite{10-kt} в~качестве основных методов можно 


выделить: 


\begin{enumerate}[(1)]


\item формальное кодирование гипотез в~качестве кортежей баз данных; 


\item оценку соответствия гипотез наборам данных;


\item построение  


при\-чин\-но-след\-ст\-вен\-ных зависимостей для параметров одной гипотезы. 


\end{enumerate}





Сущность метода \textit{формального кодирования гипотез} заключается 


в~представлении гипотез в~виде систем уравнений над кортежами базы данных. 


При этом осуществляется проверка корректности гипотез, строится схема базы 


данных и~набор функциональных зависимостей для переменных, извлеченных из 


системы уравнений. \textit{Метод оценки соответствия гипотез наборам 


данных} использует правило Байеса~\cite{11-kt}. Дополнительно в~схему базы 


данных добавляется мера соответствия гипотез исследуемому явлению. 


В~процессе оценки соответствия используется \textit{метод построения графа 


при\-чин\-но-след\-ст\-вен\-ных зависимостей} для параметров одной гипотезы.





В системе FCCE~\cite{12-kt} ключевым методом является построение попарных 


значимых корреляционных зависимостей для всех переменных множества 


разрозненных наборов данных. Для метода существует распределенная 


реализация. Метод хорошо масштабируется при сборе, извлечении и~запросах 


к~географически распределенным большим наборам данных.





В системе Robot Scientist\cite{13-kt} используются: 


\begin{enumerate}[(1)]


\item метод логического вывода 


новых гипотез при проведении экспериментов;


\item метод статической оценки 


корректности соответствия гипотез полученным данным.


\end{enumerate}





 \textit{Метод 


логического вывода} использует абдуктивное логическое 


программирование~\cite{14-kt} для генерации достоверных гипотез, 


объясняющих наблюдения, а~затем использует эти гипотезы для определения 


следующего наиболее информативного эксперимента. В~начале исследования все 


потенциальные гипотезы считаются в~равной степени верными. После этого 


устанавливается значимость оставшихся гипотез при помощи, например, теста 


Стьюдента. В~работе~\cite{15-kt} демонстрируется подход к~оценке значимости 


гипотез при помощи нескольких статистических тестов.





Основываясь на вышеприведенном анализе методов управления виртуальными 


экспериментами, можно определить \textit{основные этапы жизненного цикла 


виртуального эксперимента}.





На \textbf{первом этапе} происходит инициализация виртуального эксперимента: 


определение онтологии предметной области, определение гипотез 


и~соответству-\linebreak


\vspace*{-12pt}





\pagebreak





\noindent


ющих им моделей, определение потока работ для исполнения 


виртуального эксперимента, задание конфигурации эксперимента. Используется 


представление гипотез как специальных структур данных, сочетающих 


уравнения, входящие в~них переменные и~при\-чин\-но-след\-ст\-вен\-ные 


отображения над ними. После загрузки виртуального эксперимента в~систему 


происходит автоматическое построение решетки гипотез виртуального 


эксперимента. Решетка гипотез определяется как направленный ациклический 


граф с~вершинами, отвечающими гипотезам. Ребра графа соответствуют 


отношениям зависимости между гипотезами, когда результат вычислений одной 


гипотезы используется в~вычислениях другой гипотезы. 





На \textbf{втором этапе} 


по желанию пользователя происходит автоматическое порождение 


дополнительных гипотез из данных. 





На \textbf{третьем этапе} происходит 


построение оптимального плана исполнения эксперимента. 





На \textbf{четвертом 


этапе} происходит непосредственное проведение эксперимента, при этом 


система отслеживает, существуют ли сохраненные результаты исполнения 


отдельных функций эксперимента при заданных параметрах, и~подставляет их, 


если совпадение найде\-но. После исполнения виртуального эксперимента 


результат возвращается эксперту.





В процессе определения жизненного цикла выделены следующие \textit{операции 


в~работе над виртуальным экспериментом}. При реализации первого этапа 


к~доступным эксперту операциям над виртуальными экспериментами относятся 


до\-бав\-ле\-ние, модификация и~удаление виртуальных экспериментов, гипотез, 


потоков работ, а~также конфигураций экспериментов. Этап также включает 


автоматические операции: построение графа при\-чин\-но-след\-ст\-вен\-ных 


зависимостей между переменными одной гипотезы, специфицированной 


системой уравнений, а также построение структуры для описания взаимодействия 


между несколькими гипотезами~--- решетки гипотез. Операция построения графа 


причинно-следственных зависимостей строит граф на основании системы 


уравнений, соответствующей гипотезе. Решетка гипотез для виртуального 


эксперимента строится на основании упомянутого графа и~потока работ 


виртуального экспери-\linebreak мента.





Для реализации второго этапа используются операции автоматического 


по\-рож\-де\-ния новых гипотез из данных по желанию эксперта и~поиска 


корреляционных зависимостей между параметрами различных независимых 


гипотез.





Для реализации третьего этапа используется операция построения оптимального 


плана исполнения виртуального эксперимента. 





Для реализации четвертого этапа используется операция исполнения 


виртуального эксперимента. Промежуточные результаты расчетов, 


соответствующих вызовам отдельных функций моделей, и~оценка их 


соответствия экспериментальным данным сохраняются в~базе данных для 


последующего использования и~отслеживания эволюции гипотез и~моделей 


виртуального эксперимента. 





\section{Программная архитектура платформы исполнения 


виртуальных экспериментов}





\begin{figure}[b]


\vspace*{6pt}


 \begin{center}


 \mbox{%


 \epsfxsize=103.326mm 


 \epsfbox{kov-1.eps}


 }


\vspace*{9pt}





{\small Архитектура платформы исполнения виртуальных экспериментов}


 \end{center}


\end{figure}





В данном разделе представлена архитектура платформы (см.\ рисунок) 


и~рас\-смот\-ре\-но, каким образом компоненты архитектуры поддерживают 


различные этапы деятельности эксперта по проведению виртуальных 


экспериментов. На рисунке штриховые стрелки обозначают доступ из одного 


компонента к~интерфейсу другого компонента. Программная архитектура 


системы исполнения виртуального эксперимента включает в~себя компоненты, 


поддерживающие различные этапы проведения виртуальных экспериментов. 


В~настоящее время архитектура находится в~стадии программной реализации.





Компонент \textit{Metadata Repository} представляет собой репозиторий 


метаинформации (базу данных), предназначенный для хранения результатов 


выполнения виртуального эксперимента с~учетом накопления данных 


о~добавлении, модификации и~удалении гипотез, потоков работ, онтологий 


и~виртуальных экспериментов. Дополнительно в~репозитории сохраняются 


построенные зависимости между гипотезами (полученными либо от эксперта, 


либо автоматически извлеченные) и~имеющейся информацией о~корреляции 


между параметрами как в~рамках одной, так и~параметров нескольких гипотез.





Компонент \textit{Virtual Experiment Manager} представляет пользователю 


интерфейс с~операциями формального манипулирования виртуальными 


экспериментами и~гипотезами. После загрузки виртуального эксперимента 


в~систему его спецификация сохраняется в~базе данных. Далее по запросу 


эксперта происходит вызов компонента автоматического порождения гипотез из 


данных \textit{Hypothesis Generator}. После этого автоматически происходит 


вызов компонента \textit{Hypothesis Lattices Constructor} для построения решетки 


гипотез виртуального эксперимента. Далее вызывается компонент \textit{Virtual 


Experiment Planner}, отвечающий за построение плана исполнения виртуального 


эксперимента. Построенный план сохраняется в~базе данных. Наконец, 


происходит вызов компонента \textit{Virtual Experiment Runner}, ответственного 


за непосредственное исполнение виртуального эксперимента. Компонент 


поддерживает одновременное управление несколькими экспериментами.





Компонент \textit{Hypothesis Generator} по запросу эксперта обеспечивает 


автоматическое построение гипотез по данным в~виде нелинейных 


функциональных зависимостей. В~качестве основного метода используется 


символьная регрессия, обеспечивающая интерпретируемость полученных 


формул. Эксперту вместе с~системой уравнений возвращается оценка качества 


соответствия данным, по которой принимается решение об использовании 


построенной гипотезы. Компонент реализован в~виде модуля на языке Python 


с~использованием библиотек pandas, numpy, sympy и~deap~\cite{5-kt}. Для 


построения функциональной зависимости между переменными в~данных 


используется реализация символьной регрессии на основе генетического 


программирования с~арифметическими и~тригонометрическими 


операциями~\cite{7-kt}. 





Компонент \textit{Hypothesis Lattices Constructor} обеспечивает автоматическое 


построение решеток гипотез по соответствующим им системам уравнений. 


Компонент загружает из базы данных метаинформации спецификацию потока 


работ и~гипотез. Далее системы уравнений, соответствующих гипотезам, из 


спецификации переводятся во внутреннее представление, после чего вызывается 


компонент \textit{COA Constructor} для построения при\-чин\-но-след\-ст\-вен\-но\-го 


графа зависимостей переменных. Дальнейшая работа компонента состоит 


в~следующем. Задачи, составляющие поток работ, рассматриваются по очереди 


в~соответствии с~порядком, заданным направлением ребер потока работ. Для 


каждой пары гипотез~$I$ и~$J$, где~$I$ используется в~текущей задаче, а~$J$~--- 


в~некоторой задаче, достижимой из текущей задачи, 


конструируется транзитивное замыкание объединения структур двух гипотез 


(состоящих из уравнений и~переменных)~\cite{4-kt}. После этого производится 


объединение всех таких транзитивных замыканий. Структура результирующей 


решетки соответствует структуре объединения замыканий. Решетка гипотез 


сохраняется в~базе данных метаинформации. Компонент реализован в~виде 


модуля на языке Python с~использованием библиотек pandas, numpy и~itertools.





Компонент \textit{COA Constructor} позволяет построить причинно-следственный 


граф зависимостей переменных одной гипотезы. Компонент включает механизм 


вероятностного вывода, позволяющего отследить не только  


при\-чин\-но-след\-ст\-вен\-ные связи между гипотезами, но и~неявные 


зависимости, например значимые корреляции между параметрами гипотез, не 


связанных причинно-следственными зависимостями. Входными данными 


компонента служат пары гипотез и~решетка гипотез, а выходными~--- пары 


значимо коррелирующих переменных, соответствующих параметрам различных 


независимых гипотез. Использование пар значимо коррелирующих переменных 


позволяет системе указывать эксперту, например, на необходимость их 


совместного изменения. Компонент реализован в~виде модуля на языке Python 


с~использованием библиотек pandas, numpy и~scipy. При реализации используется 


информация о вхождении переменных в~уравнения в~виде булевой матрицы. При 


этом установление причинно-следственных связей сводится к~поиску подматриц 


с~ненулевым рангом. 





Компонент \textit{Virtual Experiment Planner} вычисляет оптимальный план 


исполнения виртуального эксперимента, при этом используются решетки гипотез, 


содержащиеся в~репозитории. Это позволяет до проведения виртуального 


эксперимента спланировать оптимальным образом подбор параметров гипотез 


и~отсечь заранее непригодные комбинации параметров. Например, уже 


вычисленные результаты вызова функции модели могут быть повторно 


использованы в~другом виртуальном эксперименте, а также могут помочь 


в~исключении виртуального эксперимента, производящего данные, не 


совпадающие с~реальными наблюдениями.





Компонент \textit{Virtual Experiment Runner} запускается для непосредственного 


исполнения виртуального эксперимента. Из репозитория метаинформации 


загружаются спецификация виртуального эксперимента и~дополнительно 


построенные решетка гипотез и~план проведения виртуального эксперимента. 


Исполнение происходит в~соответствии с~планом, при этом система использует 


сохраненные результаты исполнения части виртуального эксперимента (при их 


наличии в~репозитории метаинформации). Используя спецификацию моделей 


виртуального эксперимента, компонент осуществляет вызов соответствующих 


функций, а~затем собирает полученные данные от них. Результаты проведения 


виртуального эксперимента сохраняются в~базе данных метаинформации.





\section{Заключение}





В работе проведен анализ методов управления экспериментами, реализованных 


в~существующих системах. Рассмотрены основные этапы жизненного цикла 


виртуального эксперимента; выделены операции, используемые на каждом этапе 


жизненного цикла. Операции разделены на автоматические и~вызываемые 


экспертом. Предложена программная архитектура платформы управления 


виртуальными экспериментами, выделены ее основные компоненты, 


реализующие операции жизненного цикла виртуального эксперимента. 





Платформа находится в~стадии программной реализации и~экспериментальных 


исследований. Реализованы компоненты построения при\-чин\-но-след\-ст\-вен\-но\-го 


графа зависимостей переменных одной гипотезы, компонент автоматического 


построения решеток гипотез по соответствующим им системам уравнений, 


компонент автоматического порождения гипотез из данных, компонент 


непосредственного исполнения виртуального эксперимента. Остальные 


компоненты платформы находятся в~стадии разработки. Планируются 


экспериментальные исследования применения платформы для проведения 


виртуальных экспериментов в~нескольких областях с~интенсивным 


использованием данных, в~том числе в~астрономии и~нейрофизиологии.





{\small\frenchspacing


{%\baselineskip=10.8pt


\addcontentsline{toc}{section}{Литература}


\begin{thebibliography}{99}





\bibitem{1-kt}


\Au{Kovalev D., Kalinichenko L., Stupnikov S.} Organization of virtual experiments in  


data-intensive domains: Hypotheses and workflow specification~/\!\!/ CEUR Workshop 


Proceedings:  Selected Papers of the 19th Conference (International) on Data 


Analytics and Management in Data Intensive Domains~/


Eds. L.~Kalinichenko, Ya.~Manolopoulos, N.~Skvortsov, V.~Sukhomlin.~---


Moscow,   2017. Vol.~2022. P.~293--300.


{\sf http://ceur-ws.org/Vol-2022/paper46.pdf}.


\bibitem{2-kt}


\Au{Grau B.\,C., Horrocks I., Motik B., Parsia B., Patel-Schneider~P., 


Sattler~U.} OWL~2: The 


next step for OWL~/\!\!/ J.~Web Semant., 2008. Vol.~6. Iss.~4. P.~309--322.


\bibitem{3-kt}


\Au{Van Der Aalst W.\,M., Ter Hofstede A.\,H., Kiepuszewski B., Barros~A.\,P.} Workflow 


patterns~/\!\!/ Distrib. Parallel Dat., 2003. Vol.~14. Iss.~1. P.~5--51.


\bibitem{4-kt}


\Au{Kovalev D., Stupnikov S.} Constructing hypothesis lattices for virtual experiments in data 


intensive research~/\!\!/ Ivannikov Memorial Workshop.~--- IEEE, 2019. 


P.~14--17.


\bibitem{5-kt}


\Au{Tirikov E.} Comparison of male and female nonlinear brain functional connectivity~/\!\!/ 


CEUR Workshop Proceedings: Selected Papers of the 21st Conference 


(International) on Data Analytics and Management in Data Intensive 


Domains~/ Eds. A.~Elizarov, B.~Novikov, S.~Stupnikov.~---


Kazan,  2019. Vol.~2523. P.~404--412.


{\sf http://ceur-ws.org/Vol-2523/paper39.pdf}


\bibitem{6-kt}


\Au{Duggan J., Brodie M.} Hephaestus: Data reuse for accelerating scientific discovery~/\!\!/ 7th 


Biennial Conference on Innovative Data Systems Research: Online Proceedings, 2015. 


Paper~29. 12~p. {\sf 


http://users.eecs. northwestern.edu/$\sim$jennie/pubs/hephaestus\_full.pdf}.


\bibitem{7-kt}


\Au{Uy N.\,Q., Hoai N.\,X., O'Neill M., McKay R.\,I., 


Galv$\acute{\mbox{a}}$n-L$\acute{\mbox{o}}$pez~E.} Semantically-based 


crossover in genetic programming: Application to 


real-valued symbolic regression~/\!\!/ Genet. Program. Evol.~M., 2011. 


Vol.~12. Iss.~2. P.~91--119.


\bibitem{8-kt}


\Au{Triantafillou S., Tsamardinos I.} Constraint-based causal discovery from multiple interventions 


over overlapping variable sets~/\!\!/ J.~Mach. Learn. Res., 2015. Vol.~16. Iss.~1. 


P.~2147--2205.


\bibitem{9-kt}


\Au{Neal R.\,M.} Slice sampling~/\!\!/ Ann. Stat., 2003. 


Vol.~31. Iss.~3. P.~705--741.


\bibitem{10-kt}


\Au{Gonсalves B., Silva F., Porto F.} Upsilon-DB: A~system for data-driven hypothesis 


management and analytics~/\!\!/ arXiv.org, 2014. arXiv:1411.7419 [cs.DB]. 


6~p.


\bibitem{11-kt}


\Au{Darwiche A.} Modeling and reasoning with Bayesian networks.~--- Cambridge


University Press, 2009. 


548~p.


\bibitem{12-kt}


\Au{Schales D., Hu X., Jang J., %Sailer R., Stoecklin~M., 


\textit{et al.}} FCCE: Highly scalable 


distributed feature collection and correlation engine for low latency 


big data analytics~/\!\!/ IEEE 31st  Conference (International)


 on Data Engineering.~--- IEEE, 2015. P.~1316--1327.


\bibitem{13-kt}


\Au{King R.\,D., Whelan K.\,E., Jones F.\,M., Reiser P.\,G., Bryant~C.\,H., Muggleton~S.\,H., 


Kell~D.\,B., Oliver~S.\,G.} Functional genomic hypothesis generation and experimentation by 


a~robot scientist~/\!\!/ Nature, 2004. Vol.~427. Iss.~6971. P.~247--252.


\bibitem{14-kt}


\Au{Kakas A.\,C., Kowalski R.\,A., Toni F.} Abductive logic programming~/\!\!/ J.~Logic 


Comput., 1992. Vol.~2. Iss.~6. P.~719--770.


\bibitem{15-kt}


\Au{Тарасов Е., Ковалев Д.} Оценка качества научных гипотез в~виртуальных экспериментах 


в~областях с~интенсивным использованием данных~/\!\!/ CEUR Workshop Proceedings: 


Selected Papers of the 19th Conference (International) on Data Analytics and 


Management in Data Intensive Domains~/


Eds. L.~Kalinichenko, Ya.~Manolopoulos, N.~Skvortsov, V.~Sukhomlin.~---


Moscow,   2017. Vol.~2022. P.~281--292.


{\sf http://ceur-ws.org/Vol-2022/paper46.pdf}.





\end{thebibliography}


} }





\vspace*{-3pt}





\hfill{\small\textit{Поступила в~редакцию 15.03.20}}





\vspace*{8pt}





\hrule





\vspace*{2pt}





%\newpage





%\vspace*{-28pt}





\hrule





\vspace*{8pt}





\def\tit{ARCHITECTURE OF THE PLATFORM FOR~MANAGING HYPOTHESES-DRIVEN VIRTUAL 


EXPERIMENTS}





\def\titkol{Architecture of the platform for managing hypotheses-driven virtual 


experiments}





\def\aut{D.\,Y.~Kovalev$^1$, E.\,A.~Tarasov$^1$, V.\,N.~Zakharov$^1$, 


and~N.\,M.~Filimonov$^2$}





\def\autkol{D.\,Y.~Kovalev, E.\,A.~Tarasov, V.\,N.~Zakharov, and~N.\,M.~Filimonov}





\titel{\tit}{\aut}{\autkol}{\titkol}





\vspace*{-11pt}





\noindent


$^1$Institute of Informatics Problems, Federal Research Center ``Computer 


Science\linebreak


$\hphantom{^1}$and 


Control'' of the Russian Academy of Sciences, 44-2 Vavilov Str., Moscow\linebreak


$\hphantom{^1}$119333, 


Russian Federation





\noindent


$^2$Institute of Automatics and Computer Engineering, National Research 


University\linebreak


$\hphantom{^1}$``Moscow Power Engineering Institute,'' 14~Krasnokazarmennaya 


Str., Moscow\linebreak


$\hphantom{^1}$111250, Russian Federation 





\def\leftfootline{\small{\textbf{\thepage}


\hfill Sistemy i Sredstva Informatiki~---


Systems and Means of Informatics 2020 vol~30 no\,2}


}%


 \def\rightfootline{\small{Sistemy i Sredstva Informatiki~---


 Systems and Means of Informatics 2020 vol~30 no\,2


\hfill \textbf{\thepage}}}





\vspace*{3pt}





\Abste{The problem of carrying out virtual experiments driven by hypotheses is 


considered. The analysis of methods for managing virtual experiments in existing 


systems for working with experiments is carried out. According to the results of the 


analysis, a~life cycle of a~virtual experiment is formed. Basic operations of managing 


virtual experiments are presented for the stages of an


experiment life cycle, as well as 


the main stages of the expert's work with the platform for executing a~virtual 


experiment. The program architecture of the platform for managing virtual experiments 


and hypotheses is proposed with\linebreak


\vspace*{-12pt}}





\pagebreak





\Abstend{a~description of the main components of the platform 


and their functions that implement the basic operations of the life cycle.}





\KWE{virtual experiment life cycle; hypotheses; data intensive research}





\DOI{10.14357\!/\!08696527200206} 





%\vspace*{-24pt}





\Ack


\noindent


This research was partially supported by the Russian Foundation for Basic Research 


(project 18-07-01434~А).





%\vspace*{-12pt}





\renewcommand{\bibname}{\protect\rmfamily References}


%\renewcommand{\bibname}{\large\protect\rm References}





{\small\frenchspacing


{%\baselineskip=10.8pt


\addcontentsline{toc}{section}{References}


\begin{thebibliography}{99}





\bibitem{1-kt-1}


\Aue{Kovalev, D., L. Kalinichenko, and S. Stupnikov.} 2017. Organization of virtual 


experiments in data-intensive domains: Hypotheses and workflow specification. 


\textit{CEUR Workshop Proceedings: Selected Papers of the 19th Conference 


(International) on Data Analytics and Management in Data Intensive Domains}.


Eds. L.~Kalinichenko, Ya.~Manolopoulos, N.~Skvortsov, and V.~Sukhomlin.  Moscow.


2022:293--300. Available at:  


{\sf http://ceur-ws.org/Vol-2022/paper46.pdf} (accessed April~9, 2020).


\bibitem{2-kt-1}


\Aue{Grau, B.\,C., I. Horrocks, B. Motik, B.~Parsia, P.~Patel-Schneider, and 


U.~Sattler}. 2008. OWL 2: The next step for OWL. \textit{J.~Web Semant.} 


6(4):309--322.


\bibitem{3-kt-1}


\Aue{Van Der Aalst, W.\,M., A.\,H. Ter Hofstede, B. Kiepuszewski, and 


A.\,P.~Barros}. 2003. Workflow patterns. \textit{Distrib. Parallel Dat.} 


14(1):5--51.


\bibitem{4-kt-1}


\Aue{Kovalev, D., and S. Stupnikov.} 2019. Constructing hypothesis lattices for virtual 


experiments in data intensive research. \textit{Ivannikov Memorial Workshop}. IEEE. 14--17.


\bibitem{5-kt-1}


\Aue{Tirikov, E.} 2019. Comparison of male and female nonlinear brain functional 


connectivity. \textit{CEUR Workshop Proceedings: Selected Papers of the 21st 


Conference (International) on Data Analytics and Management in Data Intensive 


Domains.} Eds. A.~Elizarov, B.~Novikov, and S.~Stupnikov. Kazan. 2523:404--412. Available at:  


{\sf http://ceur-ws.org/Vol-2523/paper39.pdf} (accessed April~9, 2020).


\bibitem{6-kt-1}


\Aue{Duggan, J., and M. Brodie.} 2015. Hephaestus: Data reuse for accelerating 


scientific discovery. \textit{7th Biennial Conference on Innovative Data Systems 


Research: Online Proceedings.} Papre~29. 12~p. Available at: {\sf 


http://users.eecs. northwestern.edu/$\sim$jennie/pubs/hephaestus\_full.pdf} (accessed 


April~9, 2020).


\bibitem{7-kt-1}


\Aue{Uy, N.\,Q., N.\,X. Hoai, M. O'Neill, R.\,I.~McKay, and  


E.~Galv$\acute{\mbox{a}}$n-L$\acute{\mbox{o}}$pez.} 2011. Semantically-based 


crossover in genetic programming: Application to real-valued symbolic regression. 


\textit{Genet. Program. Evolv.~M.} 12(2):91--119.


\bibitem{8-kt-1}


\Aue{Triantafillou, S., and I. Tsamardinos.} 2015. Constraint-based causal discovery 


from multiple interventions over overlapping variable sets. \textit{J.~Mach. 


Learn. Res.} 16(1):2147--2205.


\bibitem{9-kt-1}


\Aue{Neal, R.\,M.} 2003. Slice sampling. \textit{Ann. Stat.} 31(3):705--741. 


\bibitem{10-kt-1}


\Aue{Gonсalves, B., F. Silva, and F. Porto.} Upsilon-DB: A~system for data-driven 


hypothesis management and analytics. \textit{arXiv.org}. 6~p. Available at: {\sf 


https://arxiv.org/ abs/1411.7419} (accessed April~9, 2020).


\bibitem{11-kt-1}


\Aue{Darwiche, A.} 2009. \textit{Modeling and reasoning with Bayesian networks}. 


Cambridge University Press. 548~p.


\bibitem{12-kt-1}


\Aue{Schales, D., X. Hu, J. Jang, %R.~Sailer, M.~Stoecklin, 


\textit{et al.}} 2015. FCCE: 


Highly scalable distributed feature collection and correlation engine for low latency big 


data analytics. \textit{31st Conference (International) on Data Engineering}. IEEE. 


1316--1327.


\bibitem{13-kt-1}


\Aue{King R.\,D., K.\,E.~Whelan, F.\,M.~Jones, P.\,G.~Reiser, C.\,H.~Bryant, 


S.\,H.~Muggleton, D.\,B.~Kell, and S.\,G.~Oliver.} 2004. Functional genomic 


hypothesis generation and experimentation by a~robot scientist. \textit{Nature} 


427(6971):247--252.


\bibitem{14-kt-1}


\Aue{Kakas, A.\,C., R.\,A. Kowalski, and F.~Toni.} 1992. Abductive logic 


programming. \textit{J.~Logic Comput.} 2(6):719--770.


\bibitem{15-kt-1}


\Aue{Tarasov E., and D. Kovalev.} 2017. Otsenka kachestva nauchnykh gipotez 


v~virtual'nykh eksperimentakh v~oblastyakh s~intensivnym ispol'zovaniem dannykh 


[Estimation of scientific hypotheses quality in virtual experiments in data intensive 


domains]. \textit{CEUR Workshop Proceedings: Selected Papers of the 19th Conference 


(International) on Data Analytics and Management in Data Intensive Domains}. 


Eds. L.~Kalinichenko, Ya.~Manolopoulos, N.~Skvortsov, and


V.~Sukhomlin. Moscow. 2022:281--292. Available at:  


{\sf http://ceur-ws.org/Vol-2022/paper46.pdf} (accessed April~9, 2020).


\end{thebibliography}


} }





%\vspace*{-3pt}





\hfill{\small\textit{Received March 15, 2020}} 





%\vspace*{-16pt} 





\Contr





\noindent


\textbf{Kovalev Dmitry Y.} (b.\ 1988)~--- scientist, Institute of Informatics Problems, 


Federal Research Center ``Computer Science and Control'' of the Russian Academy of 


Sciences, 44-2 Vavilov Str., Moscow 119333, Russian Federation; 


\mbox{dkovalev@ipiran.ru}





\vspace*{3pt}





\noindent


\textbf{Tarasov Evgeny A.} (b.\ 1987)~--- programmer, Institute of Informatics 


Problems, Federal Research Center ``Computer Science and Control'' of the Russian 


Academy of Sciences, 44-2 Vavilov Str., Moscow 119333, Russian Federation; 


\mbox{e.tarasov@outlook.com}





\vspace*{3pt}





\noindent


\textbf{Zakharov Victor N.} (b.\ 1948)~--- Doctor of Science in technology, associate 


professor, scientific secretary, Federal Research Center ``Computer Science and 


Control'' of the Russian Academy of Sciences, 44-2 Vavilov Str., Moscow 119333, 


Russian Federation; \mbox{vzakharov@ipiran.ru}





\vspace*{3pt}





\noindent


\textbf{Filimonov Nikita M.} (b.\ 1998)~--- student, Institute of Automatics and 


Computer Engineering, National Research University ``Moscow Power Engineering 


Institute,'' 14~Krasnokazarmennaya Str., Moscow 111250, Russian Federation; 


\mbox{filimonovn160@gmail.com}


\label{end\stat}





\renewcommand{\bibname}{\protect\rm Литература} 


